% Auto-split to keep each TeX source < 600 lines.

\subsubsection{Audit Kernel Hardening: KL Orthogonal Decomposition, Chernoff Breakpoint Phase Diagram, and Resonance Spectrum Ontology}\label{subsubsec:discussion-audit-kernel-hardening-kl-chernoff-resonance}
Beyond the closed loop in \S\ref{subsubsec:discussion-audit-kernel-closed-loop}, the manuscript contains three additional structural threads leaning more toward hard constraints:
they respectively objectify the intuitive notion of "orthogonal recording axes/no information leakage" into strict identities, solidify budget design into a discrete-curvature-driven breakpoint phase diagram, and pin the resonance window ontology to spectral root geometry rather than auxiliary Hankel primitive constraints.
The following statements all directly cite numbered conclusions and generated tables from the main text/appendices, introducing no new proof premises.

\paragraph{KL Pythagorean Decomposition: Strictly Externalizing Microscopic Residuals into Verifiable Fields}
At fixed resolution \(m\), the readout operator \(P_m\) and the fiber uniform lift \(Q_m\) constitute an auditable "macroscopic–microscopic" interface (\S\ref{subsec:pom-information-ledger}).
Taking the uniform distribution \(u_{\Omega_m}\) over \(\Omega_m\) as the microstate baseline, its stable pushforward is
\(
w_m:=P_m u_{\Omega_m}
\)
with \(w_m(x)=d_m(x)/2^m\) (Theorem \ref{thm:pom-kl-pythagoras-uniform}).
For any microscopic distribution \(\mu\in\Delta(\Omega_m)\), let its macroscopic histogram be \(\pi:=P_m\mu\), and let \(\mu^e:=Q_m\pi\) be the fiber-wise conditional uniform lift under the same macroscopic constraint, then the KL divergence to the uniform baseline satisfies the exact three-term identity
\[
\boxed{\;
D_{\mathrm{KL}}(\mu\|u_{\Omega_m})
=D_{\mathrm{KL}}(\pi\|w_m)\;+\;D_{\mathrm{KL}}(\mu\|\mu^e)\;}
\]
(Theorem \ref{thm:pom-kl-pythagoras-uniform}).
This strictly splits the cost of "deviation from uniform microstate" into:
\begin{itemize}
  \item \textbf{Macroscopic cost} \(D_{\mathrm{KL}}(\pi\|w_m)\): depends only on the stable-layer distribution \(\pi\) relative to the uniform-microstate pushforward \(w_m\);
  \item \textbf{Microscopic residual} \(D_{\mathrm{KL}}(\mu\|\mu^e)\): deviation of \(\mu\) relative to its fiber-wise uniformization version within the constraint class fixing \(\pi\) (average of fiber-wise non-uniformity; Lemma \ref{lem:pom-fiberwise-kl}).
\end{itemize}
Moreover, \(\mu^e\) admits simultaneous strict characterizations as both maximum entropy and I-projection:
it is the unique maximum-entropy element in the constraint set \(\mathcal{C}_\pi=\{\nu:P_m\nu=\pi\}\), equivalently the KL-nearest point to \(u_{\Omega_m}\) (Corollary \ref{cor:pom-Iproj-maxent}).
Thus the KL Pythagorean decomposition serves as the minimal verifiable definition of "orthogonal recording axes": orthogonal axes are not additional narrative coordinates, but the residual term \(D_{\mathrm{KL}}(\mu\|\mu^e)\) under the projection \(\mu\mapsto\mu^e\), whose gap characterization is also exact:
\[
D_{\mathrm{KL}}(\mu\|u_{\Omega_m})-D_{\mathrm{KL}}(\mu^e\|u_{\Omega_m})
=D_{\mathrm{KL}}(\mu\|\mu^e)
\]
(Corollary \ref{cor:pom-Iproj-gap-exact}).

\paragraph{Chernoff Quantile Certificate: Breakpoint Phase Diagram Driven by Discrete Curvature \texorpdfstring{$A_q$}{Aq}}
Writing the optimal design of tail failure rate \(\varepsilon\), resolution \(m\), and moment order \(q\) as an auditable certificate, the lower-envelope structure forces breakpoints to appear.
Specifically, for audited \(r_q\), define the secant slope \(\beta_q=(\log r_q-\log2)/(q-1)\) and certificate family
\(
\gamma_m^{(q)}(\varepsilon)=\beta_q+\frac{1}{m(q-1)}\log\frac{1}{\varepsilon}
\),
whose lower envelope
\(
\gamma_m^{\mathrm{cert}}(\varepsilon)=\inf_{q\ge 2}\gamma_m^{(q)}(\varepsilon)
\)
is a piecewise lower envelope of \(1/m\)-type hyperbolic families (Remark \ref{rem:pom-qstar-phase-diagram-breakpoints}).
Proposition \ref{prop:pom-breakpoint-curvature-law-Aq} shows: the jump of optimal moment order is completely controlled by a discrete curvature increment scalar \(A_q\),
with breakpoint closed form
\[
m_{q\to q+1}(\varepsilon)=\frac{\log(1/\varepsilon)}{A_q},
\qquad
A_q=(q-1)(\log r_{q+1}-\log2)-q(\log r_q-\log2),
\]
thus \(q^\star(m,\varepsilon)\) decreases monotonically in unit steps as \(m\) increases, with two-point simultaneous optimality at breakpoints (Remark \ref{rem:pom-qstar-phase-diagram-breakpoints}).
Within the range \(\varepsilon=10^{-6}\) and DP-extended audit \(q\le 20\), all breakpoints can be uniquely computed from \(\{r_q\}\) and tabulated in Table \ref{tab:fold_collision_qstar_breakpoints_eps1e6_q2_20};
e.g., \(q^\star=7\) for \(m\in[513.01,643.28)\), while \(q^\star=2\) for \(m\ge 5344.31\) (same remark and table).
Under the continuous interpolation regime of pressure geometry, \(A_q\) can also be interpreted as a discrete curvature sample of the analytic convex pressure surface (Remark \ref{rem:pom-breakpoint-curvature-thermo-realization}), thereby fixing the kink mechanism of the "budget design phase diagram" as a geometric consequence of curvature.

\paragraph{Resonance Window Ontology: Spectral Near-Degeneracy Rather than Hankel Primitive Degeneracy}
In the audited setting of resonance window \(q=9,\dots,17\), although Hankel-null certificates are non-empty (Table \ref{tab:fold_collision_resonance_nullmodes_q9_17}),
their "primitive quotient space dimension" is identically zero:
\[
\Delta_{\mathrm{prim}}(q)=0
\]
(Corollary \ref{cor:pom-hankel-null-no-primitive-residual}; see also Remark \ref{rem:pom-nullmodes-recurrence-plus-signature} and Table \ref{tab:fold_collision_resonance_nullmodes_primitive_quotient_q9_17} for explicit factorization decomposition).
Equivalently, any annihilating polynomial induced by a NullMode can be written as \(A^{(j)}(\lambda)=Q^{(j)}(\lambda)P_q(\lambda)\) (zero remainder), with no "primitive" null mode not explained by the minimal recurrence characteristic polynomial \(P_q\).
Thus within this window, resonance should not be understood as the appearance of "extra Hankel constraint directions", but should revert to the near-degeneracy mechanism of spectral root geometry:
the shadow spectral packet yields \(|\mu_q|/r_{q-2}\uparrow 1\) with narrow-angle window phase drift (Remarks \ref{rem:pom-shadow-spectral-packet-q9-17}, \ref{rem:pom-resonance-trimodal-phase-lock}), forming a computably one-dimensional intensity coordinate via the algebraic phase-lock residual \((\delta\theta_q,\rho_q)\) (Table \ref{tab:fold_collision_shadow_spectral_packet_algebraic_phase_residual_q9_17}; see also linearization constant and quadratic error bound: Proposition \ref{prop:pom-phase-residual-linearization-sqrt15}).
This path is conceptually consistent with the summary in Remark \ref{rem:pom-resonance-spectral-not-hankel-degeneracy}: strong resonance can manifest on the spectral side as edge-hugging slow modes and locked phases, but within the current window does not produce new Hankel primitive constraints;
thus if one wishes to "resolve resonance", one can only do so by changing the generating kernel to alter \(P_q\), or by explicitly externalizing new channels, rather than relying on excavating additional Hankel degeneracy under the same \(P_q\).

\paragraph{Minimal Field-ized Output of Audit Kernel (Hard Interface)}
The above three threads can be organized alongside the curvature error bars, Hankel normal forms, and arithmetic rigidity fields in \S\ref{subsubsec:discussion-audit-kernel-closed-loop} into a harder audit kernel output:
\begin{itemize}
  \item \textbf{Information orthogonal decomposition field:} Output \(\bigl(D_{\mathrm{KL}}(\mu\|u_{\Omega_m}),D_{\mathrm{KL}}(\pi\|w_m),D_{\mathrm{KL}}(\mu\|\mu^e)\bigr)\) enforcing the Pythagorean identity (Theorem \ref{thm:pom-kl-pythagoras-uniform}), while outputting I-projection/maximum-entropy criterion (Corollary \ref{cor:pom-Iproj-maxent}).
  \item \textbf{Moment order design field:} Taking audited \(\{r_q\}\) as input, output discrete curvature \(A_q\), breakpoints \(m_{q\to q+1}(\varepsilon)\), and phase diagram \(q^\star(m,\varepsilon)\) (Proposition \ref{prop:pom-breakpoint-curvature-law-Aq}; Remark \ref{rem:pom-qstar-phase-diagram-breakpoints}; Table \ref{tab:fold_collision_qstar_breakpoints_eps1e6_q2_20}).
  \item \textbf{Resonance spectrum field and primitivity field:} Output \(\Delta_{\mathrm{prim}}(q)\) and factorization verification (Corollary \ref{cor:pom-hankel-null-no-primitive-residual}; Table \ref{tab:fold_collision_resonance_nullmodes_primitive_quotient_q9_17}), while outputting spectral packet geometry and phase residual certificates in parallel (Remark \ref{rem:pom-shadow-spectral-packet-q9-17}; Table \ref{tab:fold_collision_shadow_spectral_packet_algebraic_phase_residual_q9_17}).
\end{itemize}
Any stable mismatch of a single field provides an implementation-neutral refutation interface, thereby rewriting the narrative descriptions of "externalization/no-leakage/budget-design/resonance-origin" as verifiable certificate consistency criteria.

